\documentclass[../../main.tex]{subfiles}
\graphicspath{{../../image/}{../../figure/}}
\begin{document}

\section{注记}\label{section:ch14-注}

前沿追踪方法由 Dafermos 于 1970 年代引入，作为构造守恒律系统 Cauchy 问题全局解的替代方法，与随机选择法并列。其基本思想是用分段线性通量逼近非线性通量，从而把 Riemann 问题的稀疏波替换为小强度稀疏激波（伪激波）的链，使得波型由有限多个激波与常状态组成。

\cref{section:ch14-标量守恒律的前沿追踪}中标量守恒律的处理是 Dafermos 原始构造的简化版本。对于单个守恒律，前沿追踪方法特别有效，因为所有激波都是可容许的，且波相互作用总是简化波型，从而保证算法在有限步内完成。

\cref{section:ch14-真非线性系统的前沿追踪}中真非线性系统的前沿追踪由 DiPerna 于 1983 年系统化，引入了近似 Riemann 解算器与简化 Riemann 解算器的两阶段策略。伪激波（以超声速速度传播的人工波）的引入是关键创新，它允许用低精度追踪小振幅波，同时以高精度追踪强度大的波。此种``混合精度''方法是控制波相互作用复杂度的核心。

\cref{section:ch14-全局波型}中波的全局身份追踪与生成阶概念由 Bressan 引入。非共振曲线与 $i$-特征的概念为后续连续依赖性与唯一性证明奠定了几何基础。

\cref{section:ch14-近似解}中 $\delta$-近似解的精确定义由 Bressan 给出。\cref{theorem:ch14-前沿追踪存在性}是前沿追踪方法的主要存在性结果，其证明在\cref{section:ch14-总变差的界}、\cref{section:ch14-伪激波总强度的界}与\cref{section:ch14-紧性与相容性}中展开。

\cref{section:ch14-总变差的界}中总变差估计与 Glimm 型泛函 $G=L+2\kappa Q$ 的引入沿用了 Glimm 随机选择法的框架。生成阶分层的思想（\cref{section:ch14-伪激波总强度的界}）由 Bressan 提出，用于证明高生成阶伪激波总强度可任意小。

\cref{section:ch14-紧性与相容性}中通过 Green 公式估计分布残差是经典的紧性--相容性论证。

\cref{section:ch14-连续依赖性}中标准 Riemann 半群的稳定性由 Bressan 于 1990 年代建立。泛函 $\rho$ 的构造是整个理论最精巧的部分：通过在激波曲线坐标系中测量两解之差，并引入权函数 $g_i=1+\kappa(Q+\bar Q)+\nu A_i$，使得 $\rho$ 在波相互作用时非增。此论证最终由 Bressan、Crandall、Grinberg 等完善。

\cref{section:ch14-标准riemann半群}中标准 Riemann 半群的抽象定义是 Bressan 的核心贡献。半群 $S_t$ 在闭子集 $\mathscr{D}\subset L^1(-\infty,\infty)$ 上定义，满足 $L^1$-Lipschitz 连续性与半群性质，其轨道即 Cauchy 问题的物理解。

\cref{section:ch14-解的唯一性}中 tame 振荡条件由 Bressan 引入，作为唯一性的充分结构条件。\cref{theorem:ch14-唯一性}的证明是前沿追踪方法最深刻的结果之一，它通过局部拼接近似解（激波扇区与线性化三角形）建立半群轨道与一般 $BV$ 解之间的相切关系。此结果由 Bressan、LeFloch 等证明。

\cref{section:ch14-连续glimm泛函与稀疏波展开}中连续 Glimm 泛函与稀疏波展开的分析由 Bressan 与 Colombo 发展，揭示了标准 Riemann 半群轨道的渐近结构：弱稀疏波按 $1/t$ 展开，最终解渐近趋于激波串。

\cref{section:ch14-强波的稳定性}中强波稳定性由 Bressan 等建立，表明强激波在小扰动下位置与强度均稳定。

前沿追踪方法与标准 Riemann 半群理论在 1990 年代完成了双曲守恒律系统弱解适定性的完整理论。后续发展包括：对非等熵气体动力学的推广、对接触间断的处理、以及对非严格双曲系统的扩展。相关文献在 Dafermos 著作的参考文献中完整列出。

\end{document}
